% Generated draft: requires mathematical and source review.
\documentclass[11pt]{article}
\usepackage[margin=25mm]{geometry}
\usepackage{fontspec}
\setmainfont{Latin Modern Roman}
\usepackage{amsmath,amssymb,mathtools,unicode-math}
\setmathfont{Latin Modern Math}
\usepackage{enumitem,needspace,xurl,hyperref}
\setlength{\parindent}{0pt}
\setlength{\parskip}{5pt}
\setlength{\emergencystretch}{3em}
\newcommand{\N}{\mathbb{N}}
\newcommand{\Z}{\mathbb{Z}}
\newcommand{\Q}{\mathbb{Q}}
\newcommand{\R}{\mathbb{R}}
\newcommand{\C}{\mathbb{C}}
\newcommand{\F}{\mathbb{F}}
\newcommand{\A}{\mathbb{A}}
\newcommand{\PP}{\mathbb{P}}
\newcommand{\E}{\mathbb{E}}
\newcommand{\eps}{\varepsilon}
\newcommand{\dd}{\,\mathrm{d}}
\begin{document}
\section*{UnsolvedMath ID: 800770}
\textbf{Source problem code:} catalogue entry
\subsection{Definitions and mathematical statement}
Let \((X,g)\) be a compact K\"ahler manifold, and let \(T_X=T_X^{1,0}\) denote its holomorphic tangent bundle. Write \(H_g(x;v)\) for the holomorphic sectional curvature in the complex tangent direction represented by \(0\ne v\in T_{X,x}\). The two curvature hypotheses in the source are
\[
H_g(x;v)\geq 0\quad\text{for every }x\in X\text{ and }0\ne v\in T_{X,x},
\]
or, separately,
\[
H_g(x;v)\leq 0\quad\text{for every }x\in X\text{ and }0\ne v\in T_{X,x}.
\]
The source introduces a subset \(V\subset T_X\) consisting of ``appropriate'' truly flat tangent vectors. It supplies neither a definition of truly flat vectors nor a rule specifying which such vectors are appropriate. Thus \(V\) is an unresolved part of the formulation; it is not assumed in advance to be a vector subbundle.

Subject to specifying this assignment \(V\), the source asks the following four questions.
\begin{enumerate}
\item Under either curvature hypothesis, does \(V\) define an integrable subbundle of \(T_X\)? In the holomorphic interpretation, this asks whether its fibers form a holomorphic vector subbundle of locally constant rank whose sheaf of local holomorphic sections is closed under Lie brackets:
\[
[\mathcal O(V),\mathcal O(V)]\subseteq\mathcal O(V).
\]
\item Under either curvature hypothesis, is there a subbundle \(W\subset T_X\) such that
\[
T_X\cong W\oplus V?
\]
If the intended meaning is an internal direct sum, the map \((w,v)\mapsto w+v\) must be an isomorphism of bundles.
\item Under the semi-positive curvature hypothesis, can one construct a maximal rationally connected (MRC) fibration
\[
f\colon X\longrightarrow Y
\]
and a complement \(W\) as above such that
\[
W=T_{X/Y}?
\]
Here an MRC fibration means a fibration whose general fibers are the equivalence classes generated by chains of rational curves through general points; the required existence and realization on \(X\) are part of the question.
\item Under the semi-negative curvature hypothesis, can one construct the Iitaka fibration associated with the canonical bundle
\[
K_X=\bigwedge^{\dim_{\C}X}T_X^*,\qquad f\colon X\longrightarrow Z,
\]
such that
\[
V=T_{X/Z}?
\]
The Iitaka fibration, when defined, is the fibration determined birationally by sufficiently divisible pluricanonical linear systems \(\lvert mK_X\rvert\). Its existence and the required realization are part of this question, rather than additional assumptions supplied in the record.
\end{enumerate}
For a holomorphic map \(f\colon X\to B\), the relative tangent sheaf is
\[
T_{X/B}=\ker\bigl(\dd f\colon T_X\to f^*T_B\bigr),
\]
with \(f^*T_B\) interpreted as the pullback tangent sheaf when necessary. This kernel is not automatically a vector subbundle. The source does not specify whether its fibration statements require everywhere defined holomorphic maps and global bundle equalities, allow meromorphic maps, or assert equalities only on suitable regular loci. It also does not specify whether the subbundles and splitting are holomorphic, smooth, or parallel for the metric connection. Consequently, the displayed questions retain unresolved qualifications and do not yet constitute a uniquely specified conjecture.
\subsection{Short English statement}
Matsumura asks whether a suitable collection of truly flat tangent vectors on a compact K\"ahler manifold with semi-positive or semi-negative holomorphic sectional curvature forms an integrable subbundle admitting a complementary subbundle. In the semi-positive case, the complement should be tangent to an MRC fibration; in the semi-negative case, the flat subbundle should be tangent to the canonical Iitaka fibration. The meaning of ``appropriate'' truly flat vectors, the bundle category, and the required global realization of the fibrations remain unspecified.
\subsection{Sources}
\begin{itemize}
\item Local UnsolvedMath registry, permanent ID \(800770\), exact source problem code \texttt{catalogue entry}, title ``Positively Curved Projective Varieties --- Problem 5,'' in the AMR Open Problem Lists collection. The supplied record attributes the problem to Matsumura, with proposed year \(2022\).
\item Supplied source reference: Matsumura, \emph{Open Problems on Structure of Positively Curved Projective Varieties} (2022), Problem 5, \url{https://arxiv.org/abs/2211.09153}. The supplied provenance records access on \texttt{2026-10-16} and extraction method \texttt{arxiv\_latex}; these are supplied provenance claims, not independent source verification.
\item Supplied registry status: \texttt{open}. The record was created and updated at \texttt{2026-10-16T14:00:00Z} and identifies upstream revision \texttt{d371ee917ac9647fbb06cd6fc529deec6f3dda82}. This status is supplied curation, not an independently verified mathematical or literature assessment.
\end{itemize}
\end{document}
